\begin{statement}{C.2.D4.3}{Proved}{Rational connected triple covers have positive fibre Euler number}
Retain the torus-fibre, rational-normalization and specified
Hirzebruch-boundary morphism of C.2.RAT.1. Suppose the horizontal
section $C$ is not contracted, its image $H$ is smooth rational,
the generic transversal along $H$ is an ordinary triple point,
and $C\to H$ has degree three. Then $e(S)\ge1$.
\end{statement}
\begin{proof}
Use the normal-cone construction of C.2.D4.1. That result forces
$N=\OO(-1)^{\oplus2}$ and a common cubic-coefficient divisor
of degree $v=1$ or $2$. The horizontal root curve is integral
of bidegree $(3,3-v)$, with rational normalization $C$.
No vertical strict conductor component dominates $H$, since
all three generic branches are supplied by $C$.

We first recall exactly how a finite attachment bound gives
the Euler inequality. In the notation of C.2.RAT.1 let $r$
be the number of original positive ruling fibres and let
$N_0$ count coefficient-zero blowup events. At least
$K\ge r-N_0$ of their positive boundary trees contain a
strict conductor curve. The negative-definite argument of
C.2.RAT.3 proves this bound without an assumption about
normalization singularities. They yield $K$ distinct
attachment points on $C$. Their images on $H$ give $k$
distinct normalization-incidence edges. Put $\ell=K-k$.
The incidence graph has first Betti number at most one,
so the exact graph count (C.RA.24) gives
$e(\Gamma_{\rm red})\ge k$. Consequently
\[
             e(S)=2-r+N_0+e(\Gamma_{\rm red})\ge2-\ell.
\tag{C.D4.13}
\]

At a point where the cubic is primitive, a simple tangent
root gives a smooth normalization branch with conductor
exactly twice the horizontal curve. Indeed the ambient
blowup chart divides the reduced equation by the third
power of its normal coordinate; the nonzero root derivative
solves the remaining coordinate holomorphically. The
remaining quadratic factor has nonzero value at the chosen
simple root. Its restriction therefore gives the conductor
factor $y^2$ times a unit, even when its own two roots coincide. This is the explicit
factorization and conductor comparison of C.2.RAT.5.
Such a point cannot carry a vertical boundary attachment.
If the horizontal root curve is locally unibranch, a
multiple root has just one point on its normalization.
A cubic has at most one multiple root unless all its roots
are simple. Thus off the common divisor each image value
receives at most one attachment when that root curve is
locally unibranch.

\paragraph{One simple common zero.}
Here the total reduced branch divisor is
$B=R_{(3,2)}+F_{(0,1)}$, with two rational normalizations
and total delta five. The complete computation C.2.RAT.6
gives $P:=\sum p_g\le1$ and
$J:=\sum b_1(Q_p)\le2$. The classes of $R,F$ are
independent. As in (C.D4.8)--(C.D4.12), writing
$\nu=\sum(r_p(B)-1)$, we obtain
\[
 b_1(U)\le2,\qquad e(U)=3\nu,\qquad
 \operatorname{rank}(\hbox{intersection pairing})
                         =2\nu+2b_1(X)-J.
\tag{C.D4.14}
\]
The torus comparison excludes $\nu\ge2$. Since $R\cdot F=3$,
we have $\nu\ge1$, so $\nu=1$. All branches of $R$ are
then locally unibranch: a multibranch point would contribute
to $\nu$ in addition to the intersection with $F$, or would
raise its contribution at that intersection to at least two.
There is exactly one point of $R\cap F$, and it has just
one preimage on $C$. Thus the common value also receives
at most one attachment. The preceding simple-root argument
handles every other value. Hence $\ell=0$ and (C.D4.13)
gives $e(S)\ge2$.

\paragraph{Two distinct common zeros.}
Now $B=R_{(3,1)}+F_1+F_2$. The curve $R$ is smooth rational,
and each $R\cap F_i$ has total intersection number three.
A contact of order $k=1,2,3$ is the plane type $A_{2k-1}$.
Let $P$ be the number of order-three contacts. The local
table gives precisely
\[
              \sum p_g=P,\qquad \sum b_1(Q_p)=2P,
                         \qquad 0\le P\le2.
\tag{C.D4.15}
\]
The cyclic cover $X_0$ has $\omega_{X_0}\simeq\OO_{X_0}$
and $\chi(\OO_{X_0})=2$, directly from its cyclic algebra
and adjunction. Proper duality identifies the quotient
$\omega_{X_0}/\rho_*\omega_X$ with a finite-length dual
module of length $P$. If $P>0$, a generator of the trivial
canonical line does not belong to the proper adjunction
ideal at at least one point. Thus no nonzero global
canonical section extends to $X$, and $p_g(X)=0$.
Leray gives $\chi(\OO_X)=2-P$. If $P=0$, the canonical
section extends and the same calculation gives $q(X)=0$.
Together these identities yield
\[
                         q(X)=\max(P-1,0).
\tag{C.D4.16}
\]
There is exactly one relation among the three strict
boundary classes: $F_1,F_2$ have the same class, while
$R$ and a fibre are independent. The exceptional matrix
is negative definite, so no further relation occurs.
Thus (C.D4.10) gives $b_1(U)=2q(X)+1\le3$. With
$\nu=\#(R\cap(F_1\cup F_2))$, the boundary topology
and Euler number give
\[
 \operatorname{rank}(\hbox{intersection pairing})
          =2\nu-4-2P+4\max(P-1,0).
\tag{C.D4.17}
\]
We check every possibility. For $P=0$ each fibre has at
least two distinct intersections, so $\nu\ge4$ and the
rank is positive, which is excluded. For $P=1$, one fibre
has a single order-three contact. The other has either
three simple intersections or one contact of order two
and one simple intersection. The first has $\nu=4$ and
positive rank; the only remaining value is $\nu=3$.
For $P=2$, each fibre has one point of intersection and
$\nu=2$. These are all the cases.

Since $R$ is smooth, the simple-root argument gives at
most one attachment at any noncommon value. A common
value has as many points on $C$ as the corresponding
intersection with $F_i$. In the remaining $P=1$ case
these numbers are one and two, so $\ell\le1$. In the
$P=2$ case they are both one, so $\ell=0$.
Equation (C.D4.13) gives $e(S)\ge1$ and $e(S)\ge2$,
respectively.

\paragraph{A double common zero.}
Finally the cubic divisor is $R_{(3,1)}+2F$. The normalized
cyclic cover has the exact eigensheaf algebra
\[
 \pi_*\OO_{X_0}
       =\OO_E\oplus L^{-1}\oplus L^{-2}(F),
       \qquad L=\OO_E(1,1).
\tag{C.D4.18}
\]
One obtains the last summand by dividing the square of the
cyclic coordinate by a local equation of the weight-two
branch; this integral element satisfies its monic cubic
relation. These summands describe the normalization in
codimension one and hence everywhere by normality. Their
$H^1$ groups vanish, since their bidegrees are $(0,0)$,
$(-1,-1)$ and $(-2,-1)$.

We verify that the normalization singularities are rational.
At a contact of order $k\le3$ the local cover is the
normalization of
\[
                            w^3=s^2(s-u^k).
\tag{C.D4.19}
\]
Blow up the common point of the two smooth branch curves
repeatedly, $k$ times, until their strict transforms are
disjoint. At each step $\sigma$, the branch multiplicity
is three, so the new cyclic line and weight-two branch
are $L'=\sigma^*L(-E)$ and $F'=\sigma^*F-E$; the exceptional
curve has branch weight zero. Both nontrivial normalized
eigensheaves on the blown-up base are therefore the
pullbacks of their predecessors tensored with $\OO(E)$.
For a point blowup,
$\sigma_*\OO(E)=\OO$ and $R^1\sigma_*\OO(E)=0$, as follows
from $0\to\OO\to\OO(E)\to\OO_E(-1)\to0$.
Finite-cover Leray then proves that the induced birational
map of the normalized covers has $R^1\OO=0$ at each step.
The final cover is smooth: each remaining branch is smooth,
the branches are disjoint, and a single branch of either
weight one or two has smooth normalization. Composing
these maps proves rationality of every singularity in
(C.D4.19). In particular $q(X)=0$, and their exceptional
configurations have $b_1=0$ (rational singularities have
rational-tree exceptional divisors, also seen here from
the indicated successive cyclic covers of $E$).

The reduced boundary has two components with independent
classes, and contains the pullback of the ample class
$(3,2)$. If $\nu=\#(R\cap F)$, the same exact homology
calculation now gives
\[
        b_1(U)=0,\qquad e(U)=3\nu,\qquad
        b_1(D)=\nu-1,\qquad
        \operatorname{rank}(\hbox{intersection pairing})=2\nu>0.
\tag{C.D4.20}
\]
This contradicts (C.D4.2), and excludes the double common
zero. The three cases exhaust $v=1,2$, proving $e(S)\ge1$.
\end{proof}


% C.2.D4.2 from part_C2_cubic_double.tex

